\section{One detector for all grids at a given height}\label{sec:detector}

Put $e(x)=\exp(2\pi i x)$, and write $G_q=(q^{-1}\mathbb Z)/\mathbb Z$
for the subgroup of order $q$ of $\mathbb R/\mathbb Z$. We construct one
test, depending on a height bound $S$ but not on $q$, to be used on every
$G_q$ with $q\le S$. All logarithms without a subscript are natural.

\begin{proposition}\label{prop:detector}
For all sufficiently large real $S$, put $t=\log S$ and
\[
 k=\left\lceil\frac{t}{\log\log t}\right\rceil,\quad
 L=\lceil t(\log t)^4\rceil,\quad
 d=\lceil\log_3(2k+1)\rceil,\qquad
 F_S(x)=\frac1L\sum_{j=1}^L e(3^{dj}x).
\]
Then $P_S=|F_S|^{2k}$ takes values in $[0,1]$, $k=o(\log S)$, and
the detector
\[
 D_S(q)=\sum_{z\in G_q}P_S(z)
 \quad\text{satisfies}\quad
 1\le D_S(q)\le q,\qquad \sum_{q\le S}D_S(q)\le S^{1+o(1)} .
\]
\end{proposition}
\begin{proof}
Expand $P_S(x)=\sum_m w_m e(mx)$. The $w_m$ form the probability
distribution of the difference of two independent sums of $k$ frequencies
drawn uniformly with replacement from $\{3^{dj}:1\le j\le L\}$.
The digits of each sum in base $3^d>k$ are its multiplicities. Thus a
collision requires the second tuple to be a permutation of the first, giving
$w_0\le k!/L^k\le (k/L)^k\le S^{-3}$ for large $S$.
Every other frequency satisfies
$\log(2+|m|)\ll dL+\log(2k)=(\log S)^{1+o(1)}$.

For any fixed $0<s<1$ and $m\ne0$, the number $b_S(m)$ of positive
divisors of $|m|$ at most $S$ satisfies
\[
 b_S(m)\le S^s \prod_{p\mid m}(1-p^{-s})^{-1}
 \le S^s\exp\bigl(O_s((\log(2+|m|))^{1-s})\bigr).
\]
Indeed there are at most $\log_2(2+|m|)$ distinct prime divisors,
the $i$th smallest is at least $i+1$, and
$-\log(1-p^{-s})\ll_s p^{-s}$. For each fixed $s$ the exponential factor
on our frequency range is $S^{o_s(1)}$, so $b_S(m)\le S^{o(1)}$
uniformly there, by taking $s$ arbitrarily small. No estimate for the
unrestricted divisor function at these frequencies is needed.
Orthogonality gives $D_S(q)=q\sum_{q\mid m}w_m$. The zero frequency
contributes at most $S^2w_0$ to the sum over $q$; each other $m$ is
charged at most $S b_S(m)$ before multiplication by $w_m$. This proves
the sum bound. The point $z=0$ contributes one and $P_S\le1$.
\end{proof}

\section{Signed sums of independent Cantor fibers}\label{sec:native}

We use $C$ also for the image of the middle-third Cantor set on the
circle. Multiplication by $3$ maps $C$ into itself. For $H=3^h$ with
$h\ge1$ write $[j]_H=1+H+\cdots+H^{j-1}$. If $Z$ is a circle set,
a centered fiber at modulus $q$ means any subset of
$(Z\cap (v+G_q))-v$, with $v$ any point of the circle. Thus independently
chosen fibers can always be added inside $G_q$. We write $b\cdot x$
or $b\cdot A$ for scalar multiplication and $bA$ for the $b$-fold
additive sumset of $A$ for positive integral $b$.

\begin{theorem}\label{thm:native}
For every positive integer $m$ and $\sigma>0$ there are $p_0,S_0$
such that if $S\ge S_0$, $p>p_0$ is prime, $H=3^h$, and
$q\le S$ divides $[p]_H$, then for any $m$ independently chosen
centered fibers $A_i$ of $(H-1)\cdot C$ at $q$,
\begin{equation}\label{eq:native-mixed}
 \left|\sum_{i=1}^m(A_i-A_i)\right|\le S^\sigma D_S(q).
\end{equation}
For fibers of $C$ itself, the same conclusion holds for all $q\le S$,
without $p$, $H$, or any divisibility condition.
\end{theorem}

The first assertion requires the same $H$ for all factors in one sum,
but no relation between their offsets or their witnesses. In particular
further restrictions on each fiber are allowed. A sum of single signed
copies of such sets, with either sign on each, is up to translation a
subset of the sum in \eqref{eq:native-mixed}: fix one element in each
nonempty set and subtract it. The proof has three steps.

\subsection*{A common frequency multiplier}
For a probability $\omega$ on a finite circle grid write
$\widehat\omega(a)=\sum_z \omega(z)e(az)$ and
$E(\omega)=(\sum_z\omega(z)^2)^{-1}$. Reflection is written
$\widetilde\omega$ and convolution is taken with counting measure.
We first allow an arbitrary fixed number $r$ of input laws $\pi_i$,
one per independently centered fiber, and set
$\lambda=\mathop{*}_{i=1}^r(\pi_i*\widetilde\pi_i)$.

Every probability on $C$ has a Fourier coefficient at some
$g\in\{1,\ldots,7\}$ of modulus at least $1/14$. Indeed the Fej\'er
kernel $\tfrac18|\sum_{a=0}^7e(a(x-1/2))|^2$ is at most
$1/(8\sin^2(\pi(x-1/2)))\le1/2$ on $C$. Its constant coefficient
is one and the absolute sum of its other coefficients is seven. Integrating
proves the assertion. It also applies to the distribution of $3^u X$
for any $C$-valued $X$ and integer $u\ge0$.

For each atom of $\pi_i$ fix a $C$-witness $X_i$ with
$(H-1)X_i-v_i$ equal to that atom modulo one, and choose the $X_i$
independently with the induced laws. Fix any sample time $u=dj$ from
Proposition~\ref{prop:detector}. For each $0\le n<M$ there are
$g_i(n)\in\{1,\ldots,7\}$ for which
\[
 \left|\prod_i \mathbb E e(g_i(n)3^u H^n X_i)\right|\ge a,
 \qquad a=14^{-r}.
\]
Pigeonhole the tuple and the angular octant of this product: on a set
$I$ of at least $M/(8\cdot7^r)$ times the tuple $(g_i)$ and octant
are fixed. Put $V_n=\prod_i e(g_i3^uH^n X_i)$ on the product probability
space. Independence and projection onto the central ray of the octant
give $|\mathbb E\sum_{n\in I}V_n|\ge a|I|/2$. Hence
\[
 \sum_{n,n'\in I}\operatorname{Re}\mathbb E(V_n\overline{V_{n'}})
 = \mathbb E\left|\sum_{n\in I}V_n\right|^2\ge a^2|I|^2/4.
\]
Take $M=M(r)$ large enough that $M/(8\cdot7^r)>8/a^2$. Removing
the diagonal, whose contribution is $|I|$, yields a pair $n>n'$
for which $|\mathbb E(V_n\overline{V_{n'}})|\ge a^2/8$. In other
words,
\begin{equation}\label{eq:native-mult}
 \prod_i |\widehat\pi_i(g_i3^u b)|\ge a^2/8,
 \qquad b=H^{n'}[n-n']_H .
\end{equation}
Here recentering contributes only constant phases. Crucially the
possibly large $b$ is the \emph{same} for every $i$. If $p>M$,
multiplication by $b$ permutes $G_q$: $H$ is prime to $[p]_H$, and
$\gcd([l]_H,[p]_H)=[\gcd(l,p)]_H=1$ for $1\le l<M$.
The identity follows by dividing
$\gcd(H^l-1,H^p-1)=H^{\gcd(l,p)}-1$ by $H-1$.
For fibers of $C$ itself one simply takes $b=1$ and applies the
seven-harmonic assertion to the genuine $3^u X_i$, before recentering;
no second moment or unit condition is needed in that case.

Pigeonhole over the $L$ sample times. One tuple $(g_i)$ and one pair
$(n,n')$ (no pair in the case of $C$) work at a fraction
$c_r>0$ of the times. Write $\mu_i=(g_i)_*\pi_i$ for pushforward
under scalar multiplication, $\nu=\mathop{*}_i
(\mu_i*\widetilde\mu_i)$ and $\nu_j=\nu^{*2^j}$ for $j\ge0$.
All Fourier coefficients of \(\nu\) and its convolution powers are nonnegative. By
\eqref{eq:native-mult}, or its $b=1$ version, for each fixed $j$
there is a constant $c_{r,j}>0$ such that
\[
 \sum_{z\in G_q} ((b)_*\nu_j)(z)F_S(z)\ge c_{r,j}.
\]
The expression on the left is real by symmetry. Jensen and
Cauchy--Schwarz give
\[
 c_{r,j}^{2k}\le \sum_z ((b)_*\nu_j)(z) P_S(z)
 \le \|\nu_j\|_2 D_S(q)^{1/2}
\]
since $P_S^2\le P_S$ and $b$ acts bijectively. Consequently
\begin{equation}\label{eq:native-scaled}
 E(\nu_j)\le D_S(q)S^{o_{r,j}(1)}.
\end{equation}
This error is uniform in the laws, $h,p,q$, and offsets once $p>M(r)$,
because $k=o(\log S)$.

\subsection*{Removing the bounded individual multipliers}
We use the weighted Balog--Szemer\'edi--Gowers lemma
\cite[Lemma~1]{TaoBSG}. In the counting-measure normalization used here,
if a symmetric probability $\rho$ on a finite abelian group satisfies
$\|\rho*\rho\|_2>K^{-1}\|\rho\|_2$, $K\ge2$, then there is a
symmetric set $B$ containing zero, with $B+B$ covered by $K^{O(1)}$
translates of $B$, $|B|\le K^{O(1)}E(\rho)$, and
$\rho(v+B)\ge K^{-O(1)}$ for some $v$. All such implied powers
are absolute. In particular $|s B|\le K^{O(s)}|B|$ for positive
integers $s$, by repeatedly using the covering.

Fix $\delta>0$, let $K=S^\delta$, and $J=\lceil2/\delta\rceil$.
All probability norms on $G_q$ lie between $S^{-1/2}$ and one.
Therefore for some $0\le j<J$ we have
$\|\nu_{j+1}\|_2>K^{-1}\|\nu_j\|_2$. Apply the lemma with
$\rho=\nu_j$. By \eqref{eq:native-scaled},
$|B|\le K^{O(1)}D_S(q)S^{o_{r,\delta}(1)}$. Each $\mu_i$ is
a convolution factor of $\nu_j$, so averaging over the other factors
gives a translate $v_i'+B$ on which $\mu_i$ has mass $\ge K^{-O(1)}$.
Thus the restriction of $\pi_i$ to
$S_i=\{x:g_i x\in v_i'+B\}$ has at least this mass.

Put $b_0=\operatorname{lcm}(1,\ldots,7)$. On the event where the
two independent samples of each $\pi_i$ belong to $S_i$, their signed
sum, multiplied by $b_0$, belongs to $(2r b_0)B$. In fact the
translate cancels for each difference and $b_0/g_i$ is a positive
integer. Multiplication by $b_0$ has kernel of size at most $b_0$ on
$G_q$. An event of probability $\ge K^{-O(r)}$ for $\lambda$ is
therefore supported on at most $b_0K^{O(r b_0)}|B|$ points. By
Cauchy--Schwarz and $b_0$ absolute, this proves
\begin{equation}\label{eq:native-effective}
 E(\lambda)\le D_S(q) S^{O_r(\delta)+o_{r,\delta}(1)}.
\end{equation}
Neither concentration of the full supports nor a factorization of their
uniform sumset law has been inferred.

\subsection*{From effective size to full sumsets}
Suppose \eqref{eq:native-mixed} fails for a fixed $m,\sigma$ with
no admissible thresholds. There is then a sequence of violations with
$S\to\infty$ and $p\to\infty$ (only $S\to\infty$ in the terminal
case). At each index keep $H,p,q,S$ fixed and let $M_r$ be the
supremum of $E(\lambda)$ over all allowed $r$-tuples of laws on
independent fibers for these same parameters. Pass to a single
subsequence on which $\log_S M_r\to e_r$ for every fixed $r$ and
$\log_S D_S(q)\to\ell$; all these exponents are in $[0,1]$.
Singleton laws exist. By \eqref{eq:native-effective}, first with
$r,\delta$ fixed and then $\delta\downarrow0$, we have
$e_*:=\sup_r e_r\le\ell$. This passage uses only fixed $r$: the
threshold $p>M(r)$ need not hold simultaneously at an individual
index for all $r$.

Choose $\gamma>0$, then one fixed $r$ with $e_r>e_*-\gamma$, and
laws giving $E(\lambda)\ge M_r/2$. Now choose a new $\delta>0$,
allowed to depend on $r$. At a plateau for the unscaled powers
$\lambda^{*2^j}$, $0\le j<J$, apply the same weighted lemma.
Since $E(\lambda^{*2^j})\le M_{r2^j}$ and there are only finitely
many possible $j$, the symmetric set $B'$ it gives has
\[
 |B'|\le S^{e_*+O(\delta)+o(1)},\qquad
 B'+B'\text{ covered by }S^{O(\delta)}\text{ translates of }B'.
\]
As before, for each $i$ restrict $\pi_i$ to an appropriate
translate of $B'$ carrying mass $\ge S^{-O(\delta)}$, normalizing
to obtain $\pi_i'$. The signed convolution $\lambda'$ of the
restricted laws is supported on $G=(2r)B'$. Pointwise domination
$\pi_i'\le S^{O(\delta)}\pi_i$ gives
$\lambda'\le S^{O(r\delta)}\lambda$ and hence
\[
 E(\lambda')\ge S^{e_*-\gamma-O(r\delta)-o(1)}.
\]

For each nonempty centered fiber $A$, take $Y\subset A$ maximal
with pairwise disjoint translates $y+G$. For the uniform law $u_Y$,
disjointness yields $E(\lambda'*u_Y)=|Y|E(\lambda')$ exactly.
Further convolution with $\widetilde u_Y$ cannot decrease effective
size, by the triangle inequality for convex combinations of translated
laws. The resulting law is allowed in $M_{r+1}$. Therefore, uniformly
for every such $A$,
\[
 |Y|\le S^{\gamma+O(r\delta)+o(1)},\qquad A\subset Y+(G-G).
\]
The inclusion is maximality. Apply it separately to the $m$ fibers of
the purported violation, using the same $G$ for them all. The total
number of choices of the $2m$ centers is at most
$S^{2m\gamma+O_m(r\delta)+o(1)}$; the remaining summands belong
to $(8mr)B'$, of size $\le S^{e_*+O_m(r\delta)+o(1)}$.
Thus the limsup of the log-base-$S$ cardinalities is at most
$e_*+2m\gamma+O_m(r\delta)$. First take $\delta\downarrow0$ for
this fixed $r$, then $\gamma\downarrow0$. The result is at most
$e_*\le\ell$, whereas the violations give a liminf at least
$\ell+\sigma$. This proves Theorem~\ref{thm:native}.
